\documentclass{beamer}
\usepackage[utf8]{inputenc}
\usepackage[T1]{fontenc}
\usepackage[french]{babel}
\usepackage{amsmath,amssymb,amsfonts}
\usepackage{tikz}
\usetikzlibrary{automata,positioning,arrows}
\usepackage{xcolor}

% Définition des couleurs des feux
\definecolor{feurouge}{RGB}{255,0,0}
\definecolor{feuorange}{RGB}{255,165,0}
\definecolor{feuvert}{RGB}{0,128,0}

% Thème de la présentation
\usetheme{Madrid}
\usecolortheme{dolphin}

\title{Modèle Mathématique pour le Système de Contrôle des Feux Tricolores}
\author{Système de Contrôle des Feux Tricolores}
\date{\today}

\begin{document}

\begin{frame}
\titlepage
\end{frame}

\begin{frame}{Table des matières}
\tableofcontents
\end{frame}

\section{Introduction}

\begin{frame}{Introduction}
\begin{itemize}
    \item Le système de contrôle des feux tricolores est un problème d'optimisation complexe
    \item Objectif: minimiser le temps d'attente global tout en assurant la sécurité
    \item Nous utiliserons plusieurs modèles mathématiques:
    \begin{itemize}
        \item Théorie des graphes pour modéliser le réseau routier
        \item Automates temporisés pour la synchronisation des feux
        \item Chaînes de Markov pour modéliser le flux de trafic
        \item Programmation linéaire pour l'optimisation des durées
    \end{itemize}
\end{itemize}
\end{frame}

\section{Modélisation du Réseau Routier}

\begin{frame}{Modélisation par Théorie des Graphes}
\begin{itemize}
    \item Le réseau routier est représenté par un graphe orienté $G = (V, E)$
    \item $V$ est l'ensemble des nœuds (intersections)
    \item $E$ est l'ensemble des arêtes (routes)
    \item Chaque arête $e \in E$ a:
    \begin{itemize}
        \item Une capacité $c(e)$ (nombre maximal de véhicules par unité de temps)
        \item Un flux $f(e)$ (nombre réel de véhicules par unité de temps)
    \end{itemize}
\end{itemize}

\begin{center}
\begin{tikzpicture}[node distance=2cm]
    \node[circle,draw] (A) {$A$};
    \node[circle,draw] (B) [right of=A] {$B$};
    \node[circle,draw] (C) [below of=A] {$C$};
    \node[circle,draw] (D) [right of=C] {$D$};
    
    \draw[->,thick] (A) -- (B) node[midway,above] {$f_{AB}$};
    \draw[->,thick] (A) -- (C) node[midway,left] {$f_{AC}$};
    \draw[->,thick] (B) -- (D) node[midway,right] {$f_{BD}$};
    \draw[->,thick] (C) -- (D) node[midway,below] {$f_{CD}$};
\end{tikzpicture}
\end{center}
\end{frame}

\begin{frame}{Contraintes de Flux}
Pour chaque nœud $v \in V$, la loi de conservation du flux s'applique:
\begin{equation}
\sum_{e \in \text{in}(v)} f(e) = \sum_{e \in \text{out}(v)} f(e)
\end{equation}

où:
\begin{itemize}
    \item $\text{in}(v)$ est l'ensemble des arêtes entrantes au nœud $v$
    \item $\text{out}(v)$ est l'ensemble des arêtes sortantes du nœud $v$
\end{itemize}

Pour chaque arête $e \in E$, la contrainte de capacité s'applique:
\begin{equation}
0 \leq f(e) \leq c(e)
\end{equation}
\end{frame}

\section{Modélisation des Feux Tricolores}

\begin{frame}{Automates Temporisés}
\begin{itemize}
    \item Chaque feu tricolore est modélisé comme un automate temporisé
    \item Un automate temporisé est un tuple $A = (L, l_0, C, E)$ où:
    \begin{itemize}
        \item $L$ est un ensemble fini d'états (Rouge, Orange, Vert)
        \item $l_0 \in L$ est l'état initial (Rouge)
        \item $C$ est un ensemble d'horloges
        \item $E \subseteq L \times \Phi(C) \times 2^C \times L$ est l'ensemble des transitions
    \end{itemize}
\end{itemize}

\begin{center}
\begin{tikzpicture}[auto,node distance=2cm,>=latex]
    \node[state,fill=feurouge!30] (R) {Rouge};
    \node[state,fill=feuvert!30] (V) [right of=R] {Vert};
    \node[state,fill=feuorange!30] (O) [below right of=R] {Orange};
    
    \draw[->] (R) edge node {$x \geq T_R$} (V);
    \draw[->] (V) edge node {$x \geq T_V$} (O);
    \draw[->] (O) edge node {$x \geq T_O$} (R);
\end{tikzpicture}
\end{center}
\end{frame}

\begin{frame}{Synchronisation des Feux}
\begin{itemize}
    \item Pour un carrefour à $n$ feux, nous avons $n$ automates temporisés
    \item La synchronisation est assurée par des contraintes temporelles
    \item Pour deux feux $i$ et $j$ qui contrôlent des flux conflictuels:
    \begin{equation}
    \text{Vert}_i \cap \text{Vert}_j = \emptyset
    \end{equation}
    \item Temps de sécurité entre le passage au rouge d'un feu et au vert d'un autre:
    \begin{equation}
    \forall i,j: \text{fin}(\text{Vert}_i) + \delta_{sécurité} \leq \text{début}(\text{Vert}_j)
    \end{equation}
\end{itemize}
\end{frame}

\section{Modélisation du Flux de Trafic}

\begin{frame}{Chaînes de Markov pour le Flux de Trafic}
\begin{itemize}
    \item Le flux de trafic est modélisé comme un processus stochastique
    \item Soit $X_t$ le nombre de véhicules à un carrefour au temps $t$
    \item La probabilité de transition est donnée par:
    \begin{equation}
    P(X_{t+1} = j | X_t = i) = p_{ij}
    \end{equation}
    \item La matrice de transition $P = [p_{ij}]$ dépend de:
    \begin{itemize}
        \item L'état des feux tricolores
        \item Les taux d'arrivée des véhicules
        \item Les taux de départ des véhicules
    \end{itemize}
\end{itemize}
\end{frame}

\begin{frame}{Équations du Flux de Trafic}
Pour un feu au vert, le taux de départ des véhicules est:
\begin{equation}
\mu_{\text{vert}} = \min(X_t, c_{\text{max}})
\end{equation}

Pour un feu au rouge ou à l'orange, le taux de départ est:
\begin{equation}
\mu_{\text{rouge/orange}} = 0
\end{equation}

Le taux d'arrivée des véhicules suit une distribution de Poisson:
\begin{equation}
P(N_t = k) = \frac{(\lambda t)^k e^{-\lambda t}}{k!}
\end{equation}

où $\lambda$ est le taux moyen d'arrivée des véhicules.
\end{frame}

\section{Optimisation des Durées}

\begin{frame}{Formulation du Problème d'Optimisation}
\begin{itemize}
    \item Objectif: minimiser le temps d'attente moyen des véhicules
    \item Variables de décision: durées des phases $T_R$, $T_V$, $T_O$ pour chaque feu
    \item Fonction objectif:
    \begin{equation}
    \min \sum_{v \in V} \sum_{t=1}^{T} w_v(t)
    \end{equation}
    où $w_v(t)$ est le temps d'attente au nœud $v$ au temps $t$
\end{itemize}
\end{frame}

\begin{frame}{Contraintes d'Optimisation}
\begin{itemize}
    \item Contraintes de durée minimale pour chaque phase:
    \begin{align}
    T_R &\geq T_{R,\min} \\
    T_V &\geq T_{V,\min} \\
    T_O &\geq T_{O,\min}
    \end{align}
    
    \item Contrainte de cycle total:
    \begin{equation}
    T_R + T_V + T_O = T_{\text{cycle}}
    \end{equation}
    
    \item Contraintes de sécurité entre feux conflictuels
\end{itemize}
\end{frame}

\section{Implémentation dans le Système}

\begin{frame}{Implémentation du Modèle dans le Code}
Dans notre implémentation FastAPI:
\begin{itemize}
    \item Les carrefours sont représentés par la classe \texttt{Carrefour}
    \item Les feux sont représentés par la classe \texttt{FeuTricolore}
    \item Le cycle des feux est géré par la fonction \texttt{cycle\_automatique}
    \item Les durées des phases sont configurables via l'interface
    \item Le modèle temporel est implémenté avec:
    \begin{equation}
    \text{temps\_cycle} = (t + \text{offset}) \mod T_{\text{cycle}}
    \end{equation}
\end{itemize}
\end{frame}

\begin{frame}{Algorithme de Synchronisation}
\begin{itemize}
    \item Pour chaque carrefour $c \in C$:
    \begin{itemize}
        \item Pour chaque feu $f \in F_c$ avec index $i$:
        \begin{itemize}
            \item Calculer le décalage: $\text{offset} = i \cdot \frac{T_{\text{cycle}}}{|F_c|}$
            \item Calculer la position dans le cycle: $t_{\text{cycle}} = (t + \text{offset}) \mod T_{\text{cycle}}$
            \item Déterminer l'état du feu selon $t_{\text{cycle}}$:
            \begin{itemize}
                \item Si $t_{\text{cycle}} < T_V$: Vert
                \item Si $T_V \leq t_{\text{cycle}} < T_V + T_O$: Orange
                \item Sinon: Rouge
            \end{itemize}
        \end{itemize}
    \end{itemize}
\end{itemize}
\end{frame}

\section{Extensions et Améliorations}

\begin{frame}{Extensions Possibles du Modèle}
\begin{itemize}
    \item Intégration de capteurs de trafic en temps réel:
    \begin{equation}
    T_V = f(\lambda_{\text{arrivée}}, \text{longueur de file})
    \end{equation}
    
    \item Algorithmes d'apprentissage par renforcement:
    \begin{equation}
    Q(s,a) \leftarrow Q(s,a) + \alpha [r + \gamma \max_{a'} Q(s',a') - Q(s,a)]
    \end{equation}
    
    \item Coordination entre carrefours adjacents:
    \begin{equation}
    \text{offset}_{i,j} = \frac{d_{i,j}}{v_{\text{moyenne}}}
    \end{equation}
    où $d_{i,j}$ est la distance entre les carrefours $i$ et $j$
\end{itemize}
\end{frame}

\begin{frame}{Conclusion}
\begin{itemize}
    \item Le modèle mathématique présenté permet:
    \begin{itemize}
        \item Une représentation formelle du système de feux tricolores
        \item Une optimisation des durées des phases
        \item Une synchronisation efficace des feux
    \end{itemize}
    \item L'implémentation actuelle utilise une version simplifiée du modèle
    \item Des extensions sont possibles pour améliorer l'efficacité du système
\end{itemize}
\end{frame}

\end{document}
